\par\medskip
\begin{center}
\textls[80]{b) Die Reihe \(F(\alpha,\beta,\gamma,x)\) bei \textsc{Euler} und \textsc{Gauss}. Die Begriffe Reihe}\\
\textls[80]{und Konvergenz.}
\end{center}
\noindent\hspace*{\parindent}Im II. Bande der \textit{Institutiones calculi integralis}, Sectio I. Cap. VIII, Problema
123\footnote{\textsc{Leonhardi Euleri} Opera, ser. I, vol. XII, 1913, S. 182.} behandelt \textsc{Euler} die Differentialgleichung
\begingroup\makeatletter\tagsleft@true\let\veqno\leqno\makeatother
\begin{equation*}
x^2(a+bx^n)\,ddy+x(c+ex^n)\,dy\,dx+(f+gx^n)y\,dx\,dx=0,\tag{3}
\end{equation*}\endgroup
\noindent indem er ihre Integration durch unendliche Reihen, die nach steigenden oder
fallenden Potenzen von \(x\) fortschreiten, gibt. Die allgemeinere Differentialgleichung,
die aus (3) hervorgeht, wenn auf der rechten Seite \(X\,dx^2\) steht,
bildet — wie schon oben angegeben wurde — den Gegenstand eines Abschnitts
von J. Fr. \textsc{Pfaffs} \textit{Disquisitiones analyticae} (1797). Im XII. Bande
der Nova Acta Petropolitana (erschienen 1801) befindet sich (S. 58—70) eine
am 3. Sept. 1778 eingereichte Abhandlung \textsc{Eulers}, \textit{Specimen transformationis
singularis serierum}, wo es heisst: »Ich betrachte die Reihe:
\begingroup\makeatletter\tagsleft@true\let\veqno\leqno\makeatother
\begin{equation*}
s=1+\frac{a\mathbin{.}b}{1\mathbin{.}c}x
+\Pi\frac{(a+1)(b+1)}{2(c+1)}x^2
+\Pi\frac{(a+2)(b+2)}{3(c+2)}x^3+\ldots,\tag{4}
\end{equation*}\endgroup
\noindent wo in üblicher Weise \(\Pi\) den Koeffizienten des vorhergehenden Gliedes bezeichnet«.
Auf S. 60 wird gezeigt, dass diese Reihe \(s\) die Differentialgleichung
\begingroup\makeatletter\tagsleft@true\let\veqno\leqno\makeatother
\begin{equation*}
0=x(1-x)\,dds+(c-(a+b+1)x)\,ds\,dx-abs\,dx\,dx\tag{5}
\end{equation*}\endgroup
\noindent befriedigt. \textsc{Euler} macht dann die Substitution
\begingroup\makeatletter\tagsleft@true\let\veqno\leqno\makeatother
\begin{equation*}
z=s(1-x)^{a+b-c}\tag{6}
\end{equation*}\endgroup
\noindent und beweist den Satz, dass \(z\) einer analogen Differentialgleichung genügt, wo
an Stelle von \(a,b,c\) die Grössen \(c-a,\ c-b,\ c\) getreten sind.\par
\clearpage\SourceLeaf{222}\ArithmeticHead{137}{R}{VIIb. VORGESCHICHTE DER DISQUISITIONES.}
\indent Dass \textsc{Gauss} seine allgemeinen Untersuchungen über die später mit
\(F(\alpha,\beta,\gamma,x)\) bezeichnete Reihe (2) an diese Abhandlung \textsc{Eulers} angeknüpft
hat, geht daraus hervor, dass die ältesten Aufzeichnungen über diesen Gegenstand,
die im Handbuch 18, Bd zerstreut zwischen astronomischen Rechnungen
stehen und aus dem Jahre 1806 stammen (vergl. die Bemerkungen Werke
X, 1, S. 330 zu dem Abdruck dieser Notizen, Abschnitt I, ebenda, S. 326—329),
unmittelbare Beziehungen zu der \textsc{Eulerschen} Abhandlung zeigen. Wir
finden nämlich im Art. 1 (a. a. O. S. 326) die in der \textit{Theoria Motus}, Art. 90\footnote{Werke VII, S. 116, vergl. auch S. 300, 301 und \textsc{Brendel}, Werke XI, 2, Abh. 3, S. 181.}
auftretenden Reihen und Kettenbrüche, für die, ebenso wie für die in den
Artikeln 2, 3 vorkommenden, daran anschliessenden Entwicklungen, die \textsc{Eulersche}
Abhandlung den allgemeinen Gesichtspunkt liefert. Das »Theorem« des
Art. 4 (a. a. O. S. 329), das wir mit Benutzung der Bezeichnung \(F(\alpha,\beta,\gamma,x)\)
durch die Gleichung
\begingroup\makeatletter\tagsleft@true\let\veqno\leqno\makeatother
\begin{equation*}
\frac{F(a,b,c,x)}{F(c-a,c-b,c,x)}=(1-x)^{-a-b+c}\tag{7}
\end{equation*}\endgroup
\noindent wiedergeben können, ist direkt der oben erwähnte \textsc{Eulersche} Satz (§ 10 der
zitierten Abhandlung), wobei \textsc{Gauss} auch vollständig die Bezeichnung \textsc{Eulers}
beibehalten hat. Der Übergang von diesem \textsc{Eulerschen} »Theorem« zu der
unmittelbar darauf folgenden »Verwandlung« erfordert die Kenntnis der Beziehung
\begingroup\makeatletter\tagsleft@true\let\veqno\leqno\makeatother
\begin{equation*}
\begin{aligned}
F(\alpha,\beta+1,\gamma+1,x)&=F(\alpha,\beta,\gamma,x)\\
&+\frac{\alpha(\gamma-\beta)}{\gamma(\gamma+1)}xF(\alpha+1,\beta+1,\gamma+2,x),
\end{aligned}\tag{8}
\end{equation*}\endgroup
\noindent die zwar zuerst in der 1809 verfassten Notiz II (Werke X, 1, S. 332) desselben
Handbuchs vorkommt (siehe die Gl. 3) des Art. 1, a. a. O., S. 332), die
aber \textsc{Gauss} hiernach schon 1806 gekannt haben muss. Mit dieser Relation
(8), die sofort die Kettenbruchentwicklung für den Quotienten
\begingroup\makeatletter\tagsleft@true\let\veqno\leqno\makeatother
\begin{equation*}
G(\alpha,\beta,\gamma,x)=\frac{F(\alpha,\beta+1,\gamma+1,x)}{F(\alpha,\beta,\gamma,x)}\tag{9}
\end{equation*}\endgroup
\noindent liefert, besass \textsc{Gauss} den Schlüssel zur gesamten formalen Theorie der Reihe
\(F(\alpha,\beta,\gamma,x)\). Wir bemerken noch, dass die Kettenbruchentwicklung durch
Anwendung des \textsc{Lambertschen} Divisionsverfahrens\footnote{J. H. \textsc{Lambert}, Histoire de l’Académie, Berlin 1768, Mémoires, S. 265.} erhalten wird, das, wie
wir sehen werden, auch im allgemeinen Falle in den \textit{Disquisitiones} (Werke III,\par
\HistoricalFoot{X 2 Abh. 2.}{18}
\clearpage\SourceLeaf{223}\ArithmeticHead{138}{L}{L. SCHLESINGER, ÜBER GAUSS' FUNKTIONENTHEORETISCHE ARBEITEN.}
\noindent S. 134) zur Geltung kommt. Dass \textsc{Gauss} auf eine Konvergenzuntersuchung
dieser Kettenbrüche nicht eingeht, hat seinen Grund wohl darin, dass \textsc{Lambert}
für die von ihm betrachteten Fälle eine solche Untersuchung gibt. Jedenfalls
hat \textsc{Gauss} diese Kettenbrüche nicht nur formal, sondern auch quantitativ
untersucht, indem es ihm ja auf ihre Brauchbarkeit für die numerische Rechnung
ankam, wie z. B. bei den in dem Briefe an \textsc{Bessel} vom 3. September
1805 (Werke X, 1, S. 237—242, besonders S. 241) auftretenden, die zur Berechnung
der Koeffizienten in der Entwicklung von \((a^2+a'^2-2aa'\cos\varphi)^{-\frac12}\)
nach den Kosinus der Vielfachen von \(\varphi\) dienen\footnote{Man vergl. etwa die Stelle in diesem Briefe a. a. O.: »Sie können hier ohne Bedenken die roten [d. h. die in der Handschrift mit roter Tinte geschriebenen] Grössen ganz vernachlässigen, da sie kaum eine Einheit in der 7. Dezimale machen können.«}.
Etwa gleichzeitig mit
dem Abschnitt I dürfte eine in der Scheda Am (S. 46, 47) enthaltene Aufzeichnung
gemacht sein, die die Differentialgleichung in derselben Form wie
bei \textsc{Euler} (oben Gl. (5), vergl. den Art. 2 der Notiz II, Werke X, 1, S. 332,
333) und die allgemeine Reihe mit den Buchstaben \(\alpha,\beta,\gamma\) (vergl. den Art. 1
der Notiz II, a. a. O. S. 332) enthält.
\par
Verschiedene Umstände mögen \textsc{Gauss} zu der Beschäftigung mit dieser
Reihe veranlasst haben. Wir haben schon oben darauf hingewiesen, wie die
Wurzeln zu dieser Arbeit bis in die neunziger Jahre des XVIII. Jahrhunderts
zurückreichen; zu jener Zeit mag sich \textsc{Gauss} besonders dafür interessiert haben,
dass die Reihenentwicklungen der vollständigen elliptischen Integrale erster
und zweiter Gattung nach Potenzen des Moduls als spezielle Fälle dieser Reihe
erscheinen. Die Bedeutung der Reihe \(F(\alpha,\beta,\gamma,x)\) für die Astronomie tritt
sowohl in den Entwicklungskoeffizienten des Ausdrucks\footnote{\textit{Disquisitiones}, Art. 6, Werke III, S. 128, vergl. auch die Abhandlung über das agM. von 1800, Werke III, S. 370, 371, ferner \textsc{Gauss} an \textsc{Bessel}, 3. Sept. 1805, Werke X, 1, S. 237.}
\[
(a^2+b^2-2ab\cos\varphi)^{-n},
\]
als auch in den Kettenbruchentwicklungen im Art. 90 der \textit{Theoria motus}\footnote{Siehe die Angaben von \textsc{Gauss}, \textit{Disquisitiones}, Art. 14, Werke III, S. 137, \textit{Anzeige}, ebenda S. 200, ferner \textsc{Riemann}, Werke, 1892, S. 84: »Diese Anwendungen, namentlich astronomische, scheinen \textsc{Gauss} zu seinen Untersuchungen \mbox{.\ .\ .\ .}\ veranlasst zu haben.« Vergl. auch \textsc{Brendel}, Werke XI, 2, Abh. 3, S. 232.}
zu Tage. Den Einfluss von Joh. Friedr. \textsc{Pfaff} haben wir auch schon oben
erwähnt. Dass die Korrespondenz mit \textsc{Bessel} über die Fakultäten zur Beschleunigung\par
\clearpage\SourceLeaf{224}\ArithmeticHead{139}{R}{VIIb. VORGESCHICHTE DER DISQUISITIONES.}
\noindent der Ausarbeitung der \textit{Disquisitiones} beigetragen hat, ist sehr
wahrscheinlich, aber einen Einfluss auf den Inhalt hat diese Korrespondenz
\textls[100]{sicher nicht} gehabt, da ja die formalen Grundlagen der Abhandlung, wie
die Aufzeichnungen in dem Handbuch 18, Bd zeigen, schon 1806 festgelegt
waren.
\par
Wenn wir noch darauf hinweisen, dass die Definition, die \textsc{Gauss} für die
sogenannte inexplikable Funktion \(\Pi x=1\mathbin{.}2\ldots x\) gewählt hat, und die er in
dem Briefe an \textsc{Bessel} vom 21. November 1811 (Werke X, 1, S. 363) und in
der Abhandlung selbst (Art. 18, Werke III, S. 144) gibt, nämlich
\begingroup\makeatletter\tagsleft@true\let\veqno\leqno\makeatother
\begin{equation*}
\begin{aligned}
\Pi(x)&=\lim_{k=\infty}\Pi(k,x),\\
\Pi(k,x)&=\frac{1\mathbin{.}2\mathbin{.}3\ldots k}{(x+1)(x+2)(x+3)\ldots(x+k)}k^x,
\end{aligned}\tag{10}
\end{equation*}\endgroup
\noindent auf \textsc{Euler} zurückgeht\begingroup\renewcommand{\thefootnote}{205a)}\footnote{\textsc{Euler} bedient sich (Brief an \textsc{Goldbach} vom 13. Oktober 1729 \textit{Corresp. math. et phys.} S. 3, vergl. \textit{Institutiones calculi differentialis} II, Cap. XVII, S. 834) zur Definition von \(x!\) des unendlichen Produkts:
\[
\frac{1^{1-x}\mathbin{.}2^x}{1+x}\mathbin{.}\frac{2^{1-x}\mathbin{.}3^x}{2+x}\mathbin{.}\frac{3^{1-x}\mathbin{.}4^x}{3+x}\ldots;
\]
\textsc{Gauss} weist (\textit{Disquisitiones}, Art. 20, Werke III, S. 145, 146) selbst — allerdings ohne \textsc{Euler} zu nennen — auf die Identität seiner Definition von \(\Pi x\) mit der \textsc{Eulerschen} hin: »Wir definieren die Funktion \(\Pi z\) durch den Wert des Produkts
\[
\frac{1\mathbin{.}2\mathbin{.}3\ldots k\mathbin{.}k^z}{(z+1)(z+2)(z+3)\ldots(z+k)}
\]
für \(k=\infty\), oder, falls man das vorzieht, durch den Grenzwert des unendlichen Produkts
\[
\frac{1}{1+z}\mathbin{.}\frac{2^{z+1}}{1^z(2+z)}\mathbin{.}\frac{3^{z+1}}{2^z(3+z)}\mathbin{.}\frac{4^{z+1}}{3^z(4+z)}\ \text{etc.«}
\]
}\addtocounter{footnote}{-1}\endgroup, so haben wir wohl die historischen Quellen der
\textsc{Gaussschen} Untersuchung sämtlich aufgezeigt.
\par
Viel ausführlicher und inhaltreicher als die im Handbuch 18, Bd enthaltenen
Aufzeichnungen ist ein aus dem Jahre 1809 stammender Entwurf
(Werke X, 1, Abschnitt III, S. 338—353) mit dem Titel: \textit{Einiges über die Reihe}
\(1+\frac{\alpha\mathbin{.}\beta}{1\mathbin{.}\gamma}x+\ldots\), der im Handbuch 19, Be aufgezeichnet ist\footnote{Siehe das Historische in den Bemerkungen, Werke X, 1, S. 353, 354.}. Er gibt in
den Artikeln 1, 2 (schon mit der Bezeichnung \(F(\alpha,\beta,\gamma)\), aber ohne das vierte
Element \(x\)) die Beziehung zwischen verwandten Reihen, und zwar immer in
der besondern Form
\begingroup\makeatletter\tagsleft@true\let\veqno\leqno\makeatother
\begin{equation*}
v=Au+Bu',\quad v'=Cu+Du',\tag{11}
\end{equation*}\endgroup
\noindent wenn \(u,v\) verwandte Reihen sind, d. h. Reihen, in denen (siehe \textit{Anzeige}, Werke\par
\HistoricalFoot{}{18\textsuperscript{*}}
\clearpage\SourceLeaf{225}\ArithmeticHead{140}{L}{L. SCHLESINGER, ÜBER GAUSS' FUNKTIONENTHEORETISCHE ARBEITEN.}
\noindent III, S. 199) »der Wert eines der drei ersten Elemente um eine Einheit verschieden,
die Werte der drei übrigen hingegen gleich sind«, und \(u',v'\) die
Derivierten nach \(x\) bedeuten. Am Schluss des Art. 2 steht die Bemerkung
»das bisherige bekannt gemacht Comm. Rec. Soc. Gott. T. II«, d. i. also in
den \textit{Disquisitiones}. Im Art. 3 finden wir die Differentialgleichung für \(F(\alpha,\beta,\gamma,x)\)
und »für den Fall, wo \(2\gamma=\alpha+\beta+1\) ist, welcher eine besondere Aufmerksamkeit
verdient«, für zwei linear unabhängige Lösungen \(P,Q\) den Wert der
Determinante
\begingroup\makeatletter\tagsleft@true\let\veqno\leqno\makeatother
\begin{equation*}
QP'-Q'P=A(x-x^2)^{-\gamma},\tag{11\,a}
\end{equation*}\endgroup
\noindent wobei der Wert der Konstanten \(A\) ohne Beweis angegeben wird. Speziell
für \(\alpha=\beta=\tfrac12,\ \gamma=1\), was dem agM. entspricht, ist dieser Wert gleich \(\frac1\pi\),
was den »Schönen Lehrsatz« (vergl. Fussnote \textsuperscript{148)}) ergibt. Eine Bestimmung
der Konstanten \(A\) für ganz beliebige \(\alpha,\beta,\gamma\) findet sich im Art. 9 (a. a. O.
S. 352) mit Hilfe der \(\Pi\)-Funktion, die im Art. 8 (a. a. O. S. 351) durch das
\textsc{Eulersche} Integral
\begingroup\makeatletter\tagsleft@true\let\veqno\leqno\makeatother
\begin{equation*}
\Pi z=\int_0^\infty e^{-x}x^z\,dx\tag{12}
\end{equation*}\endgroup
\noindent erklärt wird. — Im Art. 4 leitet \textsc{Gauss} die \textsc{Eulersche} Gleichung (7) ab, »eine
der merkwürdigsten Relationen«, wie er sagt. Es folgt die Entwicklung des
Quotienten (9) in einen Kettenbruch und in den Artikeln 5, 6 der Beweis
für die in dem Briefe an \textsc{Bessel} vom 3. September 1805 (Werke X, 1, S. 237—242)
gegebenen Entwicklungen, wobei das agM. ins Spiel gebracht und
auch ausdrücklich genannt wird.
\par
Wir beschliessen diese der Vorgeschichte der \textit{Disquisitiones} gewidmeten
Darlegungen mit dem Hinweis auf die Bedeutung, in der \textsc{Gauss} die Worte
»Reihe« und »Konvergenz« in den \textit{Disquisitiones} gebraucht.
\par
Bei den ältern Analysten und ebenso auch noch bei \textsc{Gauss} schwankt
einmal die Bedeutung des Wortes \textls[100]{Reihe} (series) und im Zusammenhang damit
auch die des Wortes \textls[100]{Konvergenz}. \textls[100]{Series} heisst ursprünglich das,
was wir jetzt gewöhnlich »Folge« nennen; so spricht \textsc{Gauss} z. B. \textit{Disquisitiones},
Art. 16 (Werke III, S. 139) von der »Reihe \(M,M',M'',\ldots\)« und \(M+M'+M''+\ldots\)
ist dann »die Summe der Reihe, deren Glieder \(M,M',M'',\ldots\) sind«
(a. a. O. S. 141, IV). In diesem Sinne heisst eine Reihe konvergent, wenn\par
\clearpage\SourceLeaf{226}\ArithmeticHead{141}{R}{VIIb. REIHE UND KONVERGENZ IN DEN DISQUISITIONES.}
\noindent ihre Glieder nach dem Grenzwert Null hin abnehmen\footnote{In diesem Sinne auch bei der »Konvergenz der Entwicklung der Mittelpunktsgleichung« Werke X, 1, S. 420.}; so sagt \textsc{Gauss} z. B.
\textit{Disquisitiones}, Art. 3 (Werke III, S. 126): »\mbox{.\ .\ .\ .}\ wird die Reihe sicher, wenn
auch nicht gleich von Anfang an, so doch nach einem gewissen Intervall konvergent
sein, und zu einer endlichen und völlig bestimmten Summe führen«.
Es kann also eine »konvergente Reihe« auch eine »unendliche Summe« haben;
so heisst es z. B. im Art. 15 (a. a. O. S. 139): »Wir werden zeigen, und zwar
denen zu Gefallen, die die strengen Methoden der antiken Geometer begünstigen,
in aller Strenge, dass \mbox{.\ .\ .\ .}
\par
\textls[100]{drittens}, die Koeffizienten ins Unendliche abnehmen, sobald \(\alpha+\beta+\gamma-1\)
eine negative Grösse ist,
\par
\textls[100]{viertens}, die Summe unserer Reihe für \(x=1\), unbeschadet der Konvergenz
im dritten Falle, unendlich ist, sobald \(\alpha+\beta-\gamma\) eine positive Grösse oder
Null ist.«
\par
Im Gegensatz zu dieser Bedeutung der Worte series und convergentia
benutzt \textsc{Gauss} z. B. gleich in der Überschrift das Wort \textls[100]{series infinita} in
dem uns geläufigen Sinne einer »unendlichen Reihe«, und Art. 3 (Werke III,
S. 126) heisst es: »Es leuchtet hiernach ein, dass solange unsere Funktion nur
als Summe der Reihe definiert ist, die Untersuchung der Natur der Sache nach
auf die Fälle beschränkt bleiben muss, wo die Reihe konvergiert, und demgemäss
die Frage keinen Sinn hat, welches der Wert der Reihe sei, für einen
Wert von \(x\), der grösser ist als Eins«, wo also die \textls[100]{Konvergenz} ebenfalls
in dem uns geläufigen Sinne zu verstehen ist. Ganz in demselben Sinne heisst
es in der \textit{Anzeige} (Werke III, S. 198): »Hier gilt eben die Reihe selbst als
Ursprung der transzendenten Funktionen \mbox{.\ .\ .}\ Die erstere Erzeugung macht,
ihrer Natur nach, die Einschränkung auf die Fälle notwendig, wo die Reihe
konvergiert, also wo das vierte Element \(x\) positiv oder negativ den Wert 1
nicht überschreitet«, und ebenso in dem Briefe an \textsc{Bessel} (Werke X, 1, S. 363)
»Eine Reihe, die nicht immer konvergiert, wie meine obige, kann auch nur
innerhalb der Schranken, wo sie konvergiert, als Definition gelten.«
\par
Wenn so, wie gesagt, auch bei \textsc{Gauss} der Sprachgebrauch für series und
und convergentia noch schwankt, so muss jedoch hervorgehoben werden, dass
für \textsc{Gauss} über die \textls[100]{Bedeutung} einer unendlichen Reihe, sowie über die\par
\clearpage\SourceLeaf{227}\ArithmeticHead{142}{L}{L. SCHLESINGER, ÜBER GAUSS' FUNKTIONENTHEORETISCHE ARBEITEN.}
\noindent \textls[100]{Sache}, ob dieselbe einen endlichen oder unendlich grossen oder überhaupt
keinen Grenzwert besitzt, keinerlei Zweifel oder Unklarheit vorhanden ist;
seine Untersuchung ist tatsächlich — wie \textsc{Pringsheim}\begingroup\renewcommand{\thefootnote}{207a)}\footnote{Encyklopädie, Bd. I, 1, S. 79.}\addtocounter{footnote}{-1}\endgroup\ bemerkt — das erste
Beispiel einer exakten Konvergenzuntersuchung im modernen Sinne.
\par
Die Frage der Konvergenz behandelt \textsc{Gauss} auch noch in zwei Handschriften,
die sich im Nachlass (Fa) befinden und die dadurch besonders charakteristisch
sind, dass die eine aus sehr früher, die andere aus sehr später
Zeit stammt. Die eine, betitelt \textit{Grundbegriffe der Lehre von den Reihen} wurde
bereits oben im Abschnitt III d) (S. 49) besprochen, die andere ist ein in fünf verschiedenen
Fassungen vorliegender Entwurf zu einer Abhandlung: \textit{(Bestimmung
der) Convergenz der Reihen, in welche die periodischen Functionen einer veränderlichen
Grösse entwickelt werden}, und stammt, da darin die Götting. gelehrten
Anzeigen von 1831 und die Jubiläumsschrift von 1849 zitiert werden, sicher
aus später Zeit; wir kommen auf diesen Entwurf\footnote{Werke X, 1, S. 400 und 407 sind zwei von den fünf Fassungen abgedruckt, vergl. die Bemerkungen ebenda, S. 437, 438, 445.} im Abschnitt X a) zurück.
